% ─────────────────────────────────────────────────────────────────────────────
% Abstract
% Paper: "Silent Decisions Are Type Errors: Enforcing AI Accountability
%         via Linear Resource Types"
% Venue: IEEE Computer — Special Issue on AI Governance
% Deadline: 13 April 2026
% Authors: Daniel Lau Pereira Soares
% ─────────────────────────────────────────────────────────────────────────────

\begin{abstract}
LGPD (Art.~18) and the EU AI Act (Art.~86) demand auditable
accountability for AI decisions with legal consequences, but
Policy-as-Code tools treat accountability as an asynchronous runtime
property, generating post-hoc logs that a failed pipeline can simply
never write --- a \emph{silent decision}. We introduce
\emph{Bound-To-Verdict} (BTV): a \texttt{Verdict} ($V$) cannot be
constructed within Safe Rust's type boundaries without atomically
consuming an \texttt{EvidenceToken} ($E$) and a
\texttt{ComplianceToken} ($C$), $V \multimap (E \otimes C)$, following
Girard's Linear Logic~\cite{girard1987linear}. We prove the
\emph{Constitutional Enclosure Theorem}: the set of materialized
silent decisions is empty within the crate's type perimeter, under the
Safe Rust and affine-ownership assumptions of \S\ref{sec:encapsulation}.
A reference implementation validates this empirically: the type-level
discipline emits no runtime code, so overhead is exactly the
cryptographic primitives' cost (BLAKE3, HMAC-SHA256) --- except making
the record \emph{durable}, which costs one to two orders of magnitude
more, measured here rather than assumed away.
\end{abstract}
